% ── ABSTRACT (~200 words) ─────────────────────────────────────────────────────
% No mathematics required. State the two-wings principle and structural
% isomorphism finding. Accessible to \bahai scholars.

\AbdulBaha stated, ``Religion and science are the two wings upon which man's intelligence can soar into the heights, with which the human soul can progress''~\citep[p.~143]{ParisTalks}. I test the principle of the talk ``The Acceptance of the Relation between Religion and Science'' against three landmarks of modern physics: Bell's theorem~\citep{Bell1964}, relational quantum mechanics~\citep{Rovelli1996}, and the experiments recognized by the 2022 Nobel Prize in Physics~\citep{Zeilinger2023Nobel}.

My finding is that the \bahai metaphysics of \mahabbat (relational love as the creative foundation of existence) shares, on my reading, a definite logical structure with relational quantum mechanics and Whitehead's process philosophy~\citep{Whitehead1929}: existence constituted by relationship, and properties emerging through relational interaction. On that structure all three agree. The further claim, that what sustains relationship is love, belongs to the writings and, in his own vocabulary, to Whitehead; physics describes the structure but names no such force.

I argue that the \bahai writings articulate this structure decades before \citeauthor{Whitehead1929} set it out philosophically, and before experiments on Bell's inequalities showed that quantum correlations cannot be explained by local properties that particles carry with them. This priority is not claimed as proof of supernatural authority; I document it as evidence of independent convergence. On my reading, the two-wings principle holds at the level of structure, not only of metaphor, though the convergence is evidence rather than proof.
